% !TEX root = ../main.tex
\lecture{5}{2026-10-29}{Datenrecht}

\section{Datenrecht}

\subsection{3 Rechtsgebiete}
\begin{itemize}
    \item Privatrecht (zwischen private, gleiche Personen)
    \item Strafrecht 
    \item öffentliches Recht (zwischen Staat und private Parteien, Regelungen ihrer Beziehungen)
\end{itemize}
\subsubsection{Privatrecht}
\begin{itemize}
    \item Personen (Definition)
    \item Personen zueinander (Beziehungen)
    \item Personen und sachen (Sachrecht)
    \item Personen und intellektuelle Güter (Patenten, Urheber, Designrecht)
    \item Austauschverhältnisse 
    \item Personengesellschaften (einfache Gesellschaft, Kollektive)
    \item Anteilsgesellschaften (AGs, GmbHs)
    \item Wertpapiere (Wertpapierrecht)
\end{itemize}

\subsubsection{Wem werden Daten zugeordnet?}
Daten gehören niemandem im rechtlichen Sinne, sondern nur nach ihrem Ursprung zugeordnet
\begin{itemize}
    \item rechtliche Zuordnung
    \begin{itemize}
          \item Patentenrecht
          \item Sachenrecht
          \item Urheberrecht
    \end{itemize}
    \item faktische Zuordnung
    \begin{itemize}
      \item Unlauteren Wettbewerbsgesetz (UWG)
      \item Strafrecht
    \end{itemize}
\end{itemize}

\subsubsection{Sachenrecht}
Eines des stärksten Rechtsformen. 
\begin{definition}
    [Sache] sind unpersönliche, körperliche, für sich bestehende Gegenstände, die der menschlichen Herrschaft unterworfen werden können.
\end{definition}

\textbf{Sind keine Sachen aber als solche behandelt}
\begin{itemize}
    \item Tiere
    \item Energie
    \item Daten
\end{itemize}
 
\textbf{Patentgesetz Schutzvoraussetzungen}
\begin{itemize}
    \item in der Industrie anwendbar
    \item Eindeutigkeit/Neuheit
    \item Nicht-Naheliegend (gotta get into it)
\end{itemize}

\textbf{Urheberrecht Schutzvoraussetzungen}
\begin{itemize}
    \item Schutzgegenstände: intellektuelle Werke
    \item Geistige Schöpfungen
    \item Individueller Charakter (einer natürlichen Person zuzuordnen)
\end{itemize}

\textbf{Software Schutz}
In jedem Fall wird nur den Code-Skript selbst geschützt nicht das Konzept.
\begin{itemize}
    \item \textbf{Patentenrecht}: Wenn individueller Charakter vorhanden 
    \item \textbf{Urheberrecht}: 
\end{itemize}

\textbf{Faktischen Zuordnung UWG}

\begin{itemize}
    \item Schutzgegenstand: Wettbewerb
    \item Geschäftsgeheimnisse: geschützt solang geheim.
    \item Marktreifes Arbeitsergebnis: z.B Datenbank
\end{itemize}